\section{How a closure is judged}
\label{sec:judge}

\subsection{Structural properties}

These can be checked without reference data, and a closure that fails one
is wrong in a way that accuracy elsewhere does not repair.
\begin{center}
\renewcommand{\arraystretch}{1.3}
\begin{tabularx}{\textwidth}{@{}P{3.3cm}YP{4.4cm}@{}}
\toprule
Property & Statement & Where it failed above \\
\midrule
Conservation & mass, momentum and energy of the closed model are conserved
  & source-form corrections \\
Realizability & closed moments belong to a non-negative distribution; a
  stress is positive semidefinite; $|\bF|\le cE$
  & Boussinesq in strong strain; quasi-normal; diffusion at low opacity \\
Well-posedness & hyperbolic or parabolic, as the physics requires
  & Burnett; Grad far from equilibrium; unconstrained learned moments \\
Entropy & an entropy inequality, where one exists & truncated expansions \\
Invariance & Galilean, rotational, and correct tensor symmetry
  & networks acting on raw components \\
Dissipation sign & the closure removes energy where the exact one does, and
  non-dissipative terms do no work
  & negative conductivity at high $k$ \\
Causality & memory acts on the past only; the kernel is analytic in the
  upper half-plane
  & acausal fits in frequency \\
Reciprocity & Onsager symmetry of coupled fluxes & independently fitted
  thermal force and heat flux \\
\bottomrule
\end{tabularx}
\end{center}

\subsection{Limits}

A closure should reduce to the trusted result where the small parameter is
small, at the right rate. For the heat flux this means the conductivity
$\kF$, then the coefficient $K$ of the leading deficit
\eqref{eq:deficit}, then the collisionless constant. Each is a separate
test, and Section~\ref{sec:between} showed that passing the first and the
third does not imply the second.

\subsection{Accuracy, in two modes}

\begin{itemize}
  \item \textbf{A priori:} the closure term against the reference closure
    term on held-out reference states, reported pointwise, in integral
    norms, spectrally, and by regime.
  \item \textbf{A posteriori:} the closed model integrated in a solver,
    against the projected reference. Report conservation drift,
    realizability violations, stability, and long-rollout statistics, not
    only short-time pointwise error.
\end{itemize}
Inequality \eqref{eq:gronwall} is the reason both are required.

\subsection{Separating the sources of error}

A reported difference between a closed model and a reference contains
four contributions, and the benchmark requires them to be reported
separately (scientific principle 9).
\begin{enumerate}
  \item \textbf{Discretization error} of the reference and of the reduced
    solver. Measured by refinement.
  \item \textbf{Ensemble-estimation error} of any target that is a
    statistic. Measured by ensemble size and confidence intervals.
  \item \textbf{Asymptotic-limit error:} the difference between the
    reference at a finite value of the small parameter and the limiting
    model. Measured by sweeping the parameter.
  \item \textbf{Model-form error} of the closure itself.
\end{enumerate}
A learned correction can be credited only with a reduction of the fourth.
A failure to recover the expected limit is a solver, estimator or
specification problem until shown otherwise, and is not evidence for a
correction.

\subsection{A ladder of tests}

The heat-flux manuscript proposes four levels, each supporting a different
claim. A closure that passes one does not automatically pass the next.
\begin{enumerate}
  \item Fix the collision operator, the dimension, the moment convention and
    the drive, and compare $\chi(k)$ with exact formulas or proved
    intervals.
  \item Test the leading deficit coefficient as well as the limiting
    conductivity.
  \item Change the collision model at matched conductivity.
  \item Run an embedded initial-value or driven experiment, which tests the
    resolved evolution.
\end{enumerate}
For learned closures add a fifth: hold out the physics, not samples. An
operator-transfer test must assign the collision-rate family, the ion
charge, or the coupling regime to the test set before any normalization or
fitting.
